% abstract.tex
\begin{abstract}
Multi-turn jailbreaks spread adversarial intent across innocent-looking conversation turns, evading safety classifiers designed for single-turn threats: Llama Guard~3 catches only 1.8\%. What makes a detector generalize to attack strategies never seen during training? We introduce DisPeD, a controlled framework where seven encoder variants differ solely in pretraining objective while sharing the same frozen classifier, isolating how domain exposure, supervision type, and label semantics each contribute to out-of-distribution robustness. Classification-based domain adaptation emerges as the only universally effective mechanism: it recovers near-perfect detection where vanilla encoders collapse (F1 0.97--1.0 vs.\ 0.26--0.31) and remains robust under mixed-LLM training where unsupervised pretraining fails entirely. Notably, scrambled labels match correct ones, revealing that the classification objective's structural role, not label content, drives generalization. Yet these advantages reverse on human-authored data, where unsupervised pretraining leads, exposing condition-specific limits. The resulting 184M-parameter pipeline achieves 24$\times$ speedup over LLM-as-judge alternatives, and few-shot adaptation eliminates false-positive inflation across deployment platforms. We release our code and data at\footnote{\url{https://github.com/anonymous-nlp-2026/persuasion_turn_jailbreak_detect}}.
\end{abstract}
